% Transcription — fonds Alexandre Grothendieck, shelfmark 94, batch 2 (pages 21-23).
% Established from the facsimile archives/batches/94/batch-02.pdf (a copy of
% 94.pdf fetched through the project's relay, https://grothendieck-relay.onrender.com/source/94.pdf;
% PDF page k = archivists' page 20+k, no cover sheet), per the skill
% transcribe-grothendieck. The folder is twenty-three pages; this is the second
% and last batch (pages 1-20 are batch 1, transcribed separately and possibly
% in parallel, so the notation here was fixed from these three pages alone).
% Pass: Opus 5.5 (claude-opus-5-5), 2026-09-26 — first pass, unchecked against the pages by a human.
%
% Pages skipped: none.
%
% Pages summarised: none. There is no typed leaf in this batch.
%
% Pages transcribed from his hand: 21, 22, 23 — blue ink, fast drafting, with
% a few pencil insertions (the marginal « (additif ?) » and the
% « (lim filt.) » on p. 21, « de transition » on p. 22). The three pages read
% as one continuous run; the formulas carry, part of the connecting prose does
% not and is left \ill{} or \uncertain{}.
%
% His pagination: none on these leaves. His numbering of statements runs on
% from batch 1: Remarque 18, « Applications », Théorème 19, Corollaire 20
% (the 20 overwritten on another numeral), Remarque 21. « th. 12 », cited
% twice, is a statement of batch 1 and is not on these pages.
%
% What the batch is about. Despite the inventory's filing under « Géométrie et
% topologie combinatoire », these pages are algebraic geometry: the end of a
% run on functors of quasi-coherent modules and their representability (the
% folder's title, « Construction de faisceaux de modules », fits). Remarque 18
% proposes calling « constructible » a functor satisfying the conditions of
% th. 12 except left exactness (ii), asks for the analogous notion for
% functors Qcoh(Y) -> Qcoh(X), and lists what should follow: composites,
% Ext^i, Tor_j, R^q f_* constructible. Under « Applications », Théorème 19:
% for f : X -> S proper, flat, of finite presentation and Z affine over X, the
% Weil restriction f_*(Z) = Prod_{X/S} Z/X is representable by a scheme affine
% over S, of finite presentation (resp. of finite type) when Z/X is. Proof:
% S affine, Z = lim Z_i with Z_i affine of finite presentation, transition maps
% closed immersions; reduction to Corollaire 20 (f_* of a closed immersion of
% affine X-schemes is a closed immersion), itself reduced by limits to S
% noetherian; first case Z = V(F), F coherent, where th. 12 gives P coherent on
% S with Hom(F, f^*(M)) = Hom(P, M), whence f_*(V(F)) = V(P); general case,
% Z closed in V(F), reduced to the subfunctor « making a section s pass
% through Z », i.e. f_*(Y) for Y = s^{-1}(Z) closed in X, representable by a
% closed subscheme of S « par l'exposé Murre, th. 3 ». Remarque 21: for any
% quasi-coherent F on X there is a quasi-coherent P on S with
% f_*(V(F)) = V(P), by passage to the limit from F of finite presentation.
%
% Notation fixed once for the batch: f_*(Z) for his f_*, \prod_{X/S} Z/X for
% the product he writes with X/S beneath; \varprojlim for his « lim » over a
% left arrow; V(F) for the vector fibre (on p. 23 once V(\underline{F}), F
% underlined, kept so); \underline{\mathrm{Hom}}, \underline{\mathrm{Ext}},
% \underline{\mathrm{Tor}} as underlined on the page; \mathrm{Qcoh}.
%
% Apparatus count: 8 \ill{}, 14 \uncertain{}.
\documentclass[11pt,a4paper]{article}
% Le préambule commun à toutes les transcriptions du fonds Grothendieck.
%
% Il tient en une idée : **l'incertitude se compose, elle ne se résout pas.**
% Une transcription qui lisse un mot douteux a détruit la seule chose pour
% laquelle on la faisait — un lecteur doit pouvoir savoir, à chaque endroit,
% ce qui était sur le feuillet et ce qui a été ajouté. Les six macros qui
% suivent sont donc l'appareil critique, et non de la décoration.
%
% Les mêmes commandes sont reconnues par `scripts/render.mjs`, qui produit la
% vue de lecture du site. Une macro ajoutée ici doit l'être aussi là-bas,
% faute de quoi le rendu échouera bruyamment — ce qui est voulu : un rendu
% silencieusement faux est pire qu'un rendu absent.

\ProvidesPackage{grothendieck}[2026/08/08 Transcriptions du fonds Grothendieck]

\usepackage{amsmath,amssymb,amsthm}
\usepackage{mathrsfs}
\usepackage{stmaryrd}
% Les diagrammes commutatifs sont partout dans ce fonds : c'est la langue
% même de la géométrie algébrique de Grothendieck, et une transcription qui
% les rendrait en prose ne transcrirait rien.
\usepackage{tikz-cd}
% Français par défaut — c'est la langue des feuillets — mais l'anglais est
% chargé aussi : la traduction et le résumé sont des documents anglais, et la
% typographie française (espaces avant les deux-points) y serait une faute.
% Ces documents basculent par \selectlanguage{english} en tête de corps.
\usepackage[english,french]{babel}
\usepackage{fontspec}
\usepackage{geometry}
\geometry{a4paper,margin=25mm,marginparwidth=28mm}
\usepackage{marginnote}
\usepackage{xcolor}
% ulem, not soul. `soul`'s \st and \ul are built on a character-by-character
% scan that XeLaTeX's font machinery breaks: the document fails with an
% undefined control sequence rather than degrading. ulem does the same two jobs
% and is Unicode-engine safe. `normalem` stops it hijacking \emph, which the
% transcriptions use for Grothendieck's own emphasis.
\usepackage[normalem]{ulem}
\usepackage{hyperref}
\hypersetup{colorlinks=true,linkcolor=black,urlcolor=black,pdfborder={0 0 0}}

% --- Métadonnées du lot -----------------------------------------------------
% Reprises telles quelles par le script de rendu, qui les affiche en tête de
% la vue de lecture. Elles fixent ce que le fichier prétend couvrir : sans
% elles, une transcription flotte sans point d'attache dans un fonds de seize
% mille feuillets.

\newcommand{\folder}[1]{\gdef\grothFolder{#1}}
\newcommand{\batch}[1]{\gdef\grothBatch{#1}}
\newcommand{\leaves}[2]{\gdef\grothFirst{#1}\gdef\grothLast{#2}}
\newcommand{\foldertitle}[1]{\gdef\grothFoldertitle{#1}}
\gdef\grothFolder{?}\gdef\grothBatch{?}\gdef\grothFirst{?}\gdef\grothLast{?}\gdef\grothFoldertitle{}

% --- Le résumé --------------------------------------------------------------
%
% Il ouvre la lecture modernisée, sous le titre « Résumé », et il est composé
% en petit. Ce n'est pas un ornement : c'est le seul passage du document qui
% ne soit pas des mathématiques, et un lecteur doit pouvoir le reconnaître
% avant de l'avoir lu — puis le sauter s'il connaît déjà le sujet, ou s'y
% arrêter s'il ne le connaît pas. Le filet qui le suit marque la reprise du
% corps.

\newenvironment{resume}
  {\small\setlength{\parskip}{0.45em}}
  {\par\medskip\hrule\medskip}

% --- Mots-clés ---------------------------------------------------------------
%
% En anglais, en clôture du résumé de la lecture modernisée : le vocabulaire
% moderne sous lequel chercher ce que les feuillets construisent. Le manifeste
% du site (`npm run manifest`) les extrait pour étiqueter la cote entière —
% cette ligne est la source unique des tags, il n'y a pas de fichier à côté.
\newcommand{\keywords}[1]{\par\smallskip\noindent
  {\footnotesize\textsc{Keywords} --- \textit{#1}}\par}

% --- L'appareil critique ----------------------------------------------------

% Le numéro de feuillet, en marge. C'est l'ancre qui permet de régler un
% désaccord sur une formule en regardant *un* feuillet plutôt qu'une chemise
% de six cents. Le site s'en sert aussi pour tourner les pages du fac-similé
% quand on fait défiler la transcription.
\newcommand{\leaf}[1]{%
  \marginnote{\footnotesize\color{gray!70}\textsf{#1}}%
  \label{leaf:#1}\ignorespaces}

% Illisible : on ne devine pas. Un mot inventé qui se lit comme les autres est
% le pire résultat possible.
\newcommand{\ill}{\textcolor{red!60!black}{[\,\ldots\,]}}

% Lecture douteuse : le mot est donné, et signalé comme incertain.
\newcommand{\uncertain}[1]{\uline{#1}}

% Ajout de l'éditeur — une lettre manquante, un mot sous-entendu.
\newcommand{\add}[1]{\textcolor{blue!50!black}{[#1]}}

% Biffé par Grothendieck lui-même. Ce qu'il a rayé fait partie du document :
% une rature dit souvent où la pensée a changé de direction.
\newcommand{\struck}[1]{\sout{#1}}

% Note du transcripteur, sur l'état matériel du feuillet.
\newcommand{\note}[1]{\par\smallskip{\small\color{gray!60!black}\textit{[#1]}}\par\smallskip}

% Note marginale de Grothendieck, distincte du corps du texte.
\newcommand{\marginal}[1]{\par\smallskip\noindent\hspace{1em}%
  {\small\color{gray!60!black}#1}\par\smallskip}

% --- Titre ------------------------------------------------------------------

\AtBeginDocument{%
  \begin{center}
    {\large\bfseries Fonds Alexandre Grothendieck --- cote n\textsuperscript{o}~\grothFolder}\\[2pt]
    {\itshape\grothFoldertitle}\\[4pt]
    {\small feuillets \grothFirst--\grothLast\ (lot \grothBatch)}\\[2pt]
    {\footnotesize\color{gray!60!black}Transcription établie d'après le fac-similé numérisé
     par l'Université de Montpellier.\\
     Les lectures incertaines sont soulignées, les illisibles notées [\,\ldots\,],
     les ajouts entre crochets.}
  \end{center}
  \vspace{1em}}

\endinput


\folder{94}
\batch{2}
\pages{21}{23}
\dating{[à partir de 1965]}
\watermark{Édition de démonstration}
\foldertitle{Dualité des modules / Construction de faisceaux de modules : notes manuscrites (s.d.)}

\begin{document}

\page{21}
\emph{Remarque 18.} Il y a lieu d'appeler un foncteur, satisfaisant aux conditions du th.~12, sauf
l'exactitude à gauche (ii), un \emph{foncteur constructible}. Il faudrait
\struck{(pour $X$ sur $Y$)}, \uncertain{définir de même} la notion de
foncteurs constructibles $\mathrm{Qcoh}(Y) \to \mathrm{Qcoh}(X)$. Il y
aurait une suite : composé de foncteurs constructibles \uncertain{est item},
foncteurs $\underline{\mathrm{Ext}}^{i}$, $\underline{\mathrm{Tor}}_{j}$,
$R^{q}f_{*}$ sont constructibles \ldots\ldots
\marginal{(additif ?)}
\note{« (additif ?) », au crayon dans l'angle supérieur droit, est relié par
un trait à « foncteur ». L'exposant de $R^{q}f_{*}$ est lu $q$ sans
certitude ; le $f_{*}$ est repassé.}

\subsection*{Applications}

\emph{Théorème 19.} Soit $f \colon X \to S$ morphisme propre et de prés.\
finie et plat. Alors pour tout $Z$ \emph{affine} sur $X$,
\[
  f_{*}(Z) = \prod_{X/S} Z/X
\]
est représentable par un \uncertain{schéma} affine sur $S$. Si $Z/X$ est de
prés.\ finie (resp.\ de type fini), alors il en est de même de
$\prod_{X/S} Z/X$ sur $S$.
\note{le $X$ de $f \colon X \to S$ est écrit en surcharge ; « affine » est
souligné trois fois et repassé. « $f_{*}(Z) =$ » est ajouté dans la marge
de gauche devant le produit.}

\emph{Dém.} On peut supposer $S$ affine, \struck{soit} donc $X$ qu.-cpt
\ill{}. \struck{\ill{}} Alors $Z = \varprojlim Z_i$
(\uncertain{$\varprojlim$ filt.}), $Z_i$ affines de présentation finie sur
$X$, d'où $f_{*}(Z) = \varprojlim f_{*}(Z_i)$, et il suffit de prouver pour la
première partie que les $f_{*}(Z_i)$ sont repr.\ par schémas affines de
prés.\ finie sur $S$. Pour la deuxième partie, il suffit alors de traiter le
cas $Z$ de t.f.\ sur $X$.
\note{« donc $X$ qu.-cpt » est ajouté au-dessus de la ligne, suivi de deux
parenthèses serrées contre le bord droit, « (qu.) » et « (\ill{} par) »,
non lues. « ($\varprojlim$ filt.) », au crayon, est inséré en interligne et
renvoyé par un trait après « $Z_i$, » ; il se termine par une petite flèche
non identifiée.}

\page{22}
Mais on sait que l'on peut prendre les morphismes de transition
$Z_j \to Z_i$ des immersions fermées, et il suffit de vérifier que les
$f_{*}(Z_j) \to f_{*}(Z_i)$ sont aussi des immers.\ fermées, car à la limite
\uncertain{on aura} que $f_{*}(Z) \hookrightarrow f_{*}(Z_i)$ \uncertain{par}
imm.\ fermée, donc $f_{*}(Z)$ sera de t.f.\ sur $S$. On est donc ramené à
prouver en même temps le
\note{« de transition », au crayon, est inséré en interligne au-dessus du
début de la ligne et renvoyé devant « $Z_j \to Z_i$ » ; « le » est ajouté en
interligne au-dessus de « Corollaire ».}

\emph{Corollaire 20.} Si $Z \to Z'$ est une immers.\ fermée de $X$-schémas
affines, alors $f_{*}(Z) \to f_{*}(Z')$ est une immers.\ fermée.
\note{le numéro 20 est écrit en surcharge d'un autre chiffre ; l'énoncé est
marqué d'un trait vertical dans la marge.}

Pour prouver ceci, on est \ill{}, par \ill{} standard, ramené aux cas où
$Z, Z'$ de présentation finie.

Ayant des situations de présentation finie, on est ramené au cas $S$ de type
fini sur $\mathrm{Spec}\,\mathbb{Z}$, donc noeth.

\emph{1\textsuperscript{er} Cas.} $Z$ est un fibré vectoriel $V(F)$, avec $F$
cohérent sur $X$. Alors le th.~12 nous fournit un $P$ cohérent sur $S$, tel
que
\[
  \underline{\mathrm{Hom}}(F, f^{*}(M)) \simeq \underline{\mathrm{Hom}}(P, M),
\]
et on voit \uncertain{tout de} suite que $f_{*}(V(F)) \simeq V(P)$.

\emph{2\textsuperscript{ème} Cas.} Cas général. On sait qu'on a

\page{23}
une immersion fermée
\[
  Z \hookrightarrow Z' = V(\underline{F}),
\]
et il suffit de prouver que
\[
  f_{*}(Z) \to f_{*}(Z')
\]
est représentable par immersions fermées. (\uncertain{En même} temps, on
aura prouvé le cor.~20 !). On est ramené \uncertain{à} la situation
\uncertain{ci-contre} :

\begin{tikzcd}
Z \arrow[r, hook] \arrow[d] & Z' \arrow[dl] \\
X \arrow[ur, bend right=40] \arrow[d, no head] & \\
S &
\end{tikzcd}

\note{le schéma est dessiné dans la marge de gauche ; la page écrit
« $Z \subset Z'$ » sur une ligne, rendu ici par une flèche d'inclusion. La
flèche courbe de $X$ vers $Z'$ est la section $s$ ; la lettre n'est pas
portée sur la flèche.}

une section $s$ de $Z'$ sur $X$, et on veut représenter le ss-foncteur de $S$
(\uncertain{en} $T/S$ qui \struck{\ill{}} « fait passer $s$ par $Z$ »). Soit
$s^{-1}(Z) = Y \subset X$, sous-schéma fermé de $X$. \struck{On} Le foncteur
considéré n'est autre que
\[
  f_{*}(Y) = \prod_{X/S} Y/X .
\]
On sait qu'il est représentable par un sous-schéma fermé de $S$ par l'exposé
\uncertain{Murre}, th.~3.
\note{« de $S$ » est ajouté en interligne au-dessus de « ss-foncteur ».}

cqfd.

\emph{Remarque 21.} Soit $f \colon X \to S$ propre, plat, de prés.\ finie.
Pour tout Module quasi-coh.\ $F$ sur $X$, on peut trouver un Module
quasi-coh.\ $P$ sur $S$ tel que
\[
  f_{*}(V(F)) \simeq V(P) :
\]
on procède par passage à la limite à partir du cas $F$ de prés.\ finie.
\note{les deux $V$ de cette ligne sont repassés à l'encre plus foncée.}

\end{document}
